%----------------------------------------------------------------------------------------
% GELIJKVORMIGE DRIEHOEKEN. Alle opgaven zijn eigen werk.
%
% Opbouw: Probleem, Start, leerstof, en op een nieuwe pagina de Opdrachten en Test jezelf (geen Koppel in de 2de graad).
% Labels van dit hoofdstuk beginnen met h2-.
% Pakketten: de bewijzen van ZZZ en ZHZ staan alleen in pakket 5; de hoogtelijn in een rechthoekige
% driehoek en de link met een negatieve factor alleen in pakket 4 en 5. Alleen in pakket 5 (wenken bij LPD 9 en
% het servicedocument voor 5 uur): de stelling van Thales met de evenwijdige projectie, middenparallel en
% zwaartelijnen, en de metrische betrekkingen.
%----------------------------------------------------------------------------------------

\mychapter{2}{Gelijkvormige driehoeken}

\begin{probleem}
\vspace{-1mm}\naast{Van een schoendoos maak je een eenvoudige camera. Prik vooraan een klein gaatje en span achteraan, 30~cm verder, een vel calqueerpapier. Richt de doos op een kerktoren die 150~m verder staat. Op het papier verschijnt de toren, op zijn kop en 7~cm hoog. Hoe hoog is de toren? Je kunt een verhouding opstellen, maar waarom mag dat eigenlijk?}{\begin{tikzpicture}[scale=.62,font=\scriptsize]
  \fill[accent!25] (-.35,0) rectangle (.35,2.6);
  \fill[accent!25] (-.35,2.6) -- (0,3.5) -- (.35,2.6) -- cycle;
  \draw[accent,thick] (-.35,0) -- (-.35,2.6) -- (0,3.5) -- (.35,2.6) -- (.35,0);
  \draw[lijngrijs] (-.8,0) -- (6.2,0);
  \draw[grijs,thick] (5,-1) rectangle (6.2,.7);
  \draw[accent] (0,3.5) -- (6.2,-.7);
  \draw[accent,very thick] (6.2,0) -- (6.2,-.7);
  \fill[white] (5,0) circle (2.2pt); \fill (5,0) circle (1.2pt);
  \node[below] at (2.5,0) {150 m};
  \node[below right] at (6.2,-.25) {7 cm};
  \node[right] at (.35,1.7) {?};
  \draw[grijs,|-|] (5,-1.3) -- (6.2,-1.3) node[midway,below] {30 cm};
\end{tikzpicture}}
\end{probleem}

\lkr{Leerplan \pakketnaam: LPD 9 (gelijkvormigheidskenmerken van driehoeken toepassen om meetkundige problemen op te lossen). De Verken laat het kenmerk HH eerst onderzoeken, zoals de wenk bij dat doel voorstelt. Het bewijs van HH steunt op de homothetie uit hoofdstuk 1 en hoort bij \alleenvoor{1,2,3,4}{LPD 4 (redeneringen beargumenteren; de bewijzen van de kenmerken reconstrueren is daar een extra)}\alleenvoor{5}{LPD B1 (bewijzen; het leerplan noemt de kenmerken en de stelling van Thales als voorbeelden), en daarom krijgen ook ZZZ en ZHZ een kort bewijs}. \alleenvoor{1,2,3,4}{De stelling van Thales (de evenwijdige met een zijde) is een extra bij LPD 9, als toepassing van de kenmerken. }\alleenvoor{5}{Wenken bij LPD 9 zijn de stelling van Thales en de metrische betrekkingen in een rechthoekige driehoek; de leerlijn voor 5 uur (servicedocument) maakt ze verplicht, leidt Thales in met een evenwijdige projectie, en vraagt met de kenmerken de eigenschappen van de middenparallel en de zwaartelijnen te bewijzen. Dat staat in de onderdelen over Thales, middenparallel en zwaartelijnen en de hoogtelijn. }\alleenvoor{4}{De hoogtelijn in een rechthoekige driehoek noemt het leerplan \pakketnaam\ niet: ze is een toepassing van HH, en staat er omdat hoofdstuk 3 (de middelevenredige) erop steunt. }Voorkennis: de hoekensom in een driehoek, overstaande hoeken, hoeken bij evenwijdige rechten en een snijlijn (F- en Z-hoeken), de congruentiekenmerken.}

In hoofdstuk 1 zag je dat gelijkvormige figuren dezelfde hoeken hebben, en zijden die allemaal met dezelfde factor vermenigvuldigd zijn. Bij driehoeken is er iets bijzonders aan de hand: je hoeft niet alles na te gaan. Twee gelijke hoeken volstaan al. Daarmee kun je afstanden berekenen die je niet kunt meten, zoals de hoogte van een toren of de breedte van een rivier. In het handboek Driehoeksmeting bouw je daarop verder: sinus, cosinus en tangens zijn verhoudingen in gelijkvormige driehoeken.

\begin{oefeningen}[Start]
\oef{h2-start-1}

\oef{h2-verken}
\end{oefeningen}

%========================================================================================
\section{Gelijkvormige driehoeken}

\subsection{Overeenkomstige hoekpunten}

Twee driehoeken zijn gelijkvormig als hun hoeken gelijk zijn en hun zijden dezelfde verhouding hebben. Maar welke zijde hoort bij welke? Bij een rechthoek was dat duidelijk: de korte bij de korte. Bij een driehoek die gedraaid of gespiegeld is, moet je goed kijken. Daarom schrijf je de hoekpunten die bij elkaar horen in dezelfde volgorde.

\begin{theorie}
Driehoek $ABC$ is \textbf{\textcolor{accent}{gelijkvormig}} met driehoek $DEF$, genoteerd $\formule{\drieh{ABC} \gv \drieh{DEF}}$, als

\vspace{1mm}
\qquad $\formule{\hoek{A} = \hoek{D}, \quad \hoek{B} = \hoek{E}, \quad \hoek{C} = \hoek{F}} \qquad \text{en} \qquad \formule{\dfrac{\lengte{DE}}{\lengte{AB}} = \dfrac{\lengte{EF}}{\lengte{BC}} = \dfrac{\lengte{DF}}{\lengte{AC}} = k}$

\vspace{1mm}
De volgorde van de letters zegt welke hoekpunten bij elkaar horen: $A$ bij $D$, $B$ bij $E$ en $C$ bij $F$.
\end{theorie}

Uit de volgorde lees je ook de overeenkomstige zijden af: $[AB]$ hoort bij $[DE]$, want dat zijn de eerste twee letters. En $[BC]$ hoort bij $[EF]$, de laatste twee.

\voorbeeld{}{Gegeven: $\drieh{ABC} \gv \drieh{PQR}$, met $\lengte{AB} = 4$, $\lengte{BC} = 6$, $\lengte{AC} = 5$ en $\lengte{PQ} = 6$. Bereken $\lengte{QR}$ en $\lengte{PR}$.\label{vb:h2-def}}{%
\uitleg{Uit de volgorde van de letters volgt dat $[PQ]$ bij $[AB]$ hoort. Van dat paar kennen we beide lengtes, dus daaruit volgt de factor.}
\naast{\begin{stappenvrij}{5.4cm}{4cm}
k = \dfrac{\lengte{PQ}}{\lengte{AB}} = \dfrac64 = 1{,}5 & Van $ABC$ naar $PQR$. \\[2ex]
\lengte{QR} = 1{,}5 \cdot \lengte{BC} = 1{,}5 \cdot 6 = 9 & $[QR]$ hoort bij $[BC]$. \\[1ex]
\lengte{PR} = 1{,}5 \cdot \lengte{AC} = 1{,}5 \cdot 5 = 7{,}5 & $[PR]$ hoort bij $[AC]$.
\end{stappenvrij}}{\begin{tikzpicture}[scale=.36,font=\scriptsize]
  \filldraw[fill=accent!30,draw=accent,thick] (0,0) coordinate (A) -- (4,0) coordinate (B) -- (.625,4.961) coordinate (C) -- cycle;
  \node[below] at (A) {$A$}; \node[below] at (B) {$B$}; \node[above] at (C) {$C$};
  \filldraw[fill=accent!12,draw=accent,thick] (13,0) coordinate (P) -- (7,0) coordinate (Q) -- (12.06,7.44) coordinate (R) -- cycle;
  \node[below] at (P) {$P$}; \node[below] at (Q) {$Q$}; \node[above] at (R) {$R$};
\end{tikzpicture}}
In de figuur is $PQR$ gespiegeld, maar dat maakt niets uit: de volgorde van de letters zegt wat bij elkaar hoort.}

Vaak staat de volgorde er niet bij, en zie je alleen welke hoeken gelijk zijn. Dan zoek je eerst de hoekpunten die bij elkaar horen. Een zijde hoort bij de zijde die tegenover dezelfde hoek ligt.

\voorbeeld{}{In driehoek $ABC$ is $\hoek{A} = 70^\circ$ en $\hoek{B} = 50^\circ$. In driehoek $KLM$ is $\hoek{K} = 50^\circ$ en $\hoek{M} = 70^\circ$. De driehoeken zijn gelijkvormig, met $\lengte{AB} = 8$, $\lengte{BC} = 10$ en $\lengte{MK} = 12$. Schrijf de gelijkvormigheid in de juiste volgorde en bereken $\lengte{KL}$.\label{vb:h2-volgorde}}{%
\uitleg{Gelijke hoeken horen bij elkaar: $A$ bij $M$ (beide $70^\circ$) en $B$ bij $K$ (beide $50^\circ$). Dan hoort $C$ bij $L$.}
\begin{stappen}
\drieh{ABC} \gv \drieh{MKL} & De hoekpunten in dezelfde volgorde. \\[1ex]
k = \dfrac{\lengte{MK}}{\lengte{AB}} = \dfrac{12}{8} = 1{,}5 & $[MK]$ hoort bij $[AB]$: de eerste twee letters. \\[2ex]
\lengte{KL} = 1{,}5 \cdot \lengte{BC} = 1{,}5 \cdot 10 = 15 & $[KL]$ hoort bij $[BC]$: de laatste twee letters.
\end{stappen}}

\begin{werkwijze}
Zoek welke hoekpunten bij elkaar horen: die met gelijke hoeken. Schrijf de gelijkvormigheid met de hoekpunten in dezelfde volgorde. Bereken de factor uit een paar overeenkomstige zijden waarvan je beide lengtes kent. Vermenigvuldig met $k$ om van de eerste naar de tweede driehoek te gaan, en deel door $k$ om terug te gaan.
\end{werkwijze}

\nujij{$\drieh{DEF} \gv \drieh{XYZ}$ met $\lengte{DE} = 10$, $\lengte{EF} = 8$ en $\lengte{XY} = 4$. Bereken $\lengte{YZ}$.}{$k = \tfrac{4}{10} = 0{,}4$, dus $\lengte{YZ} = 0{,}4 \cdot 8 = 3{,}2$.}

\begin{oefeningen}
\oef{h2-def}
\end{oefeningen}

%========================================================================================
\section{Kenmerken van gelijkvormige driehoeken}

\subsection{Twee gelijke hoeken}

Moet je altijd drie hoeken en drie verhoudingen nagaan om te weten dat twee driehoeken gelijkvormig zijn? In de Verken tekende je twee driehoeken met dezelfde twee hoeken, en kwamen de verhoudingen van de zijden vanzelf gelijk uit. Is dat altijd zo?

Twee gelijke hoeken geven meteen ook een gelijke derde hoek, want de som van de hoeken is altijd $180^\circ$. De vraag is dus of de zijden dan vanzelf evenredig zijn. Het idee: schuif de ene driehoek op de andere, zodat een gelijke hoek samenvalt. Dan zie je een homothetie uit hoofdstuk 1.

\begin{theorie}
\begin{minipage}[t]{.6\linewidth}
\vspace{0pt}
\begin{stelling}{Kenmerk HH}
Zijn twee hoeken van een driehoek gelijk aan twee hoeken van een andere driehoek, dan zijn de driehoeken gelijkvormig.
\end{stelling}

\begin{bewijs}
Neem $\hoek{D} = \hoek{A}$ en $\hoek{E} = \hoek{B}$. Leg $\drieh{DEF}$ op $\drieh{ABC}$, met $D$ op $A$ en $E$ op de halfrechte $[AB$. Omdat $\hoek{D} = \hoek{A}$, valt $F$ dan op $[AC$. Noem de nieuwe punten $E'$ en $F'$.

Omdat $\hoek{E'} = \hoek{B}$ (F-hoeken), is $E'F' \parallel BC$. Neem de homothetie met centrum $A$ en factor $k = \tfrac{\lengte{AE'}}{\lengte{AB}}$. Die beeldt $B$ af op $E'$, en $BC$ op de evenwijdige door $E'$. Dus $C$ gaat naar $F'$, en $\drieh{AE'F'}$ is het beeld van $\drieh{ABC}$.
\end{bewijs}
\end{minipage}\hfill
\begin{minipage}[t]{.36\linewidth}
\vspace{0pt}
\centering
\begin{tikzpicture}[scale=1.05,font=\scriptsize]
  \coordinate (A) at (0,0); \coordinate (B) at (4,0); \coordinate (C) at (1.2,2.6);
  \coordinate (E) at (2.4,0); \coordinate (F) at (.72,1.56);
  \filldraw[fill=accent!10,draw=accent,thick] (A) -- (B) -- (C) -- cycle;
  \filldraw[fill=accent!30,draw=accent,thick] (A) -- (E) -- (F) -- cycle;
  \hoekboog{B}{A}{C}{6mm}
  \hoekboog{C}{B}{A}{5mm}\hoekboog{C}{B}{A}{6mm}
  \hoekboog{F}{E}{A}{5mm}\hoekboog{F}{E}{A}{6mm}
  \node[below] at (A) {$A = D$}; \node[below] at (B) {$B$}; \node[above] at (C) {$C$};
  \node[below] at (E) {$E'$}; \node[left] at (F) {$F'$};
\end{tikzpicture}
\end{minipage}
\end{theorie}

Dit kenmerk is het belangrijkste van het hoofdstuk: twee gelijke hoeken, en je weet dat alle zijden in dezelfde verhouding staan. Je noemt het kort \emph{HH}, van hoek-hoek.

\voorbeeld{}{In driehoek $ABC$ is $\hoek{A} = 35^\circ$ en $\hoek{B} = 80^\circ$. In driehoek $DEF$ is $\hoek{D} = 65^\circ$ en $\hoek{E} = 35^\circ$. Zijn de driehoeken gelijkvormig? Schrijf de gelijkvormigheid in de juiste volgorde.\label{vb:h2-hh}}{%
\uitleg{We berekenen de derde hoek in beide driehoeken, en zoeken dan gelijke hoeken.}
\naast{\begin{stappenvrij}{4.6cm}{3.4cm}
\hoek{C} = 180^\circ - 35^\circ - 80^\circ = 65^\circ & De hoekensom in $\drieh{ABC}$. \\[1ex]
\hoek{F} = 180^\circ - 65^\circ - 35^\circ = 80^\circ & De hoekensom in $\drieh{DEF}$. \\[1ex]
\hoek{A} = \hoek{E}, \ \hoek{B} = \hoek{F} & Twee paren gelijke hoeken: kenmerk HH. \\[1ex]
\drieh{ABC} \gv \drieh{EFD} & $A$ bij $E$, $B$ bij $F$, dus $C$ bij $D$.
\end{stappenvrij}}{\begin{tikzpicture}[scale=.75,font=\scriptsize]
  \coordinate (A) at (0,0); \coordinate (B) at (3,0); \coordinate (C) at (2.67,1.87);
  \filldraw[fill=accent!30,draw=accent,thick] (A) -- (B) -- (C) -- cycle;
  \hoekboog{B}{A}{C}{4mm} \hoekboog{C}{B}{A}{2.5mm}\hoekboog{C}{B}{A}{3.1mm}
  \node[below] at (A) {$A$}; \node[below] at (B) {$B$}; \node[above] at (C) {$C$};
  \begin{scope}[shift={(4.6,3.4)},rotate=-90]
  \coordinate (E) at (0,0); \coordinate (F) at (3.6,0); \coordinate (D) at (3.2,2.24);
  \filldraw[fill=accent!12,draw=accent,thick] (E) -- (F) -- (D) -- cycle;
  \hoekboog{F}{E}{D}{4mm} \hoekboog{D}{F}{E}{2.5mm}\hoekboog{D}{F}{E}{3.1mm}
  \end{scope}
  \node[above] at (E) {$E$}; \node[below] at (F) {$F$}; \node[right] at (D) {$D$};
\end{tikzpicture}}}

\begin{oefeningen}
\oef{h2-hh}
\end{oefeningen}

\subsection{Drie zijden in dezelfde verhouding}

Soms ken je geen hoeken, maar wel de zijden: je hebt ze gemeten, of ze staan op een plan. Zijn de drie verhoudingen gelijk, dan zijn de driehoeken gelijkvormig. Dat klopt omdat de drie zijden een driehoek helemaal vastleggen: dat weet je van het congruentiekenmerk ZZZ.

\begin{theorie}
\begin{stelling}{Kenmerk ZZZ}
Hebben de drie zijden van een driehoek dezelfde verhouding tot de drie zijden van een andere driehoek, dan zijn de driehoeken gelijkvormig:

\vspace{1mm}
\qquad $\formule{\dfrac{\lengte{DE}}{\lengte{AB}} = \dfrac{\lengte{EF}}{\lengte{BC}} = \dfrac{\lengte{DF}}{\lengte{AC}} \ \Rightarrow \ \drieh{ABC} \gv \drieh{DEF}}$
\end{stelling}
\begin{alleenvoor}{5}

\begin{bewijs}
Noem de gemeenschappelijke verhouding $k$, en vergroot $\drieh{ABC}$ met factor $k$. Het beeld heeft zijden $k\,\lengte{AB} = \lengte{DE}$, $k\,\lengte{BC} = \lengte{EF}$ en $k\,\lengte{AC} = \lengte{DF}$. Het beeld en $\drieh{DEF}$ hebben dus dezelfde drie zijden, en zijn congruent (ZZZ). Ze hebben dezelfde hoeken, en $\drieh{ABC}$ heeft die hoeken ook, want een homothetie behoudt hoeken.
\end{bewijs}
\end{alleenvoor}
\end{theorie}

Welke zijde hoort bij welke? Bij gelijkvormige driehoeken hoort de kortste zijde bij de kortste, en de langste bij de langste. Rangschik de zijden dus eerst van klein naar groot.

\voorbeeld{}{Zijn een driehoek met zijden 4, 6 en 7 en een driehoek met zijden 15, 17,5 en 10 gelijkvormig?\label{vb:h2-zzz}}{%
\uitleg{We rangschikken de zijden van de tweede driehoek: 10, 15 en 17,5. Dan delen we telkens de overeenkomstige zijden.}
\begin{stappen}
\dfrac{10}{4} = 2{,}5, \quad \dfrac{15}{6} = 2{,}5, \quad \dfrac{17{,}5}{7} = 2{,}5 & Kortste bij kortste, middelste bij middelste, langste bij langste. \\[2ex]
\text{gelijkvormig, } k = 2{,}5 & Kenmerk ZZZ.
\end{stappen}}

Eén verhouding die verschilt, is genoeg om te weten dat de driehoeken niet gelijkvormig zijn.

\voorbeeld{}{Zijn een driehoek met zijden 4, 6 en 7 en een driehoek met zijden 8, 12 en 15 gelijkvormig?\label{vb:h2-zzzniet}}{%
\uitleg{Beide zijn al gerangschikt. We delen telkens.}
\begin{stappen}
\dfrac84 = 2, \quad \dfrac{12}{6} = 2, \quad \dfrac{15}{7} \approx 2{,}14 & De laatste verhouding verschilt. \\[2ex]
\text{niet gelijkvormig} & De tweede driehoek is in verhouding iets langer.
\end{stappen}}

\begin{oefeningen}
\oef{h2-zzz}
\end{oefeningen}

\subsection{Twee zijden en de hoek ertussen}

Er is nog een derde kenmerk, dat lijkt op het congruentiekenmerk ZHZ. Ken je twee zijden en de hoek ertussen, dan ligt een driehoek vast. Zijn twee zijden evenredig en is de hoek ertussen gelijk, dan zijn de driehoeken gelijkvormig.

\begin{theorie}
\begin{minipage}[t]{.62\linewidth}
\vspace{0pt}
\begin{stelling}{Kenmerk ZHZ}
Hebben twee zijden van een driehoek dezelfde verhouding tot twee zijden van een andere driehoek, en is de ingesloten hoek gelijk, dan zijn de driehoeken gelijkvormig:

\vspace{1mm}
\qquad $\formule{\dfrac{\lengte{DE}}{\lengte{AB}} = \dfrac{\lengte{DF}}{\lengte{AC}} \text{ en } \hoek{D} = \hoek{A} \ \Rightarrow \ \drieh{ABC} \gv \drieh{DEF}}$
\end{stelling}
\begin{alleenvoor}{5}

\begin{bewijs}
Noem de verhouding $k$, en leg $\drieh{DEF}$ op $\drieh{ABC}$ met $D$ op $A$. Omdat $\hoek{D} = \hoek{A}$, valt $E$ op $[AB$ en $F$ op $[AC$. De homothetie met centrum $A$ en factor $k$ beeldt $B$ af op $E$, want $\lengte{AE} = k\,\lengte{AB}$, en zo ook $C$ op $F$. Dus is $\drieh{DEF}$ het beeld van $\drieh{ABC}$.
\end{bewijs}
\end{alleenvoor}
\end{minipage}\hfill
\begin{minipage}[t]{.34\linewidth}
\vspace{0pt}
\centering
\begin{tikzpicture}[scale=.8,font=\scriptsize]
  \coordinate (A) at (0,0); \coordinate (B) at (2,0); \coordinate (C) at (.9,1.5);
  \filldraw[fill=accent!30,draw=accent,thick] (A) -- (B) -- (C) -- cycle;
  \hoekboog{B}{A}{C}{3.5mm}
  \node[below] at (A) {$A$}; \node[below] at (B) {$B$}; \node[above] at (C) {$C$};
  \begin{scope}[xshift=2.5cm]
  \coordinate (D) at (0,0); \coordinate (E) at (3,0); \coordinate (F) at (1.35,2.25);
  \filldraw[fill=accent!12,draw=accent,thick] (D) -- (E) -- (F) -- cycle;
  \hoekboog{E}{D}{F}{3.5mm}
  \end{scope}
  \node[below] at (D) {$D$}; \node[below] at (E) {$E$}; \node[above] at (F) {$F$};
\end{tikzpicture}
\end{minipage}
\end{theorie}

Let op: de gelijke hoek moet tussen de twee zijden liggen. Ligt hij er niet tussen, dan kun je soms twee verschillende driehoeken tekenen, en zegt het kenmerk niets.

\voorbeeld{}{In driehoek $ABC$ is $\lengte{AB} = 6$, $\lengte{AC} = 9$ en $\hoek{A} = 40^\circ$. In driehoek $PQR$ is $\lengte{PQ} = 4$, $\lengte{PR} = 6$ en $\hoek{P} = 40^\circ$. Zijn de driehoeken gelijkvormig? Welke hoek van $PQR$ is even groot als $\hoek{B}$?\label{vb:h2-zhz}}{%
\uitleg{De hoek van $40^\circ$ ligt in beide driehoeken tussen de twee gegeven zijden. We vergelijken de verhoudingen, de kortste zijde met de kortste.}
\begin{stappen}
\dfrac{\lengte{AB}}{\lengte{PQ}} = \dfrac64 = 1{,}5, \quad \dfrac{\lengte{AC}}{\lengte{PR}} = \dfrac96 = 1{,}5 & Dezelfde verhouding. \\[2ex]
\hoek{A} = \hoek{P} = 40^\circ & De ingesloten hoeken zijn gelijk. \\[1ex]
\drieh{ABC} \gv \drieh{PQR} & Kenmerk ZHZ, met $A$ bij $P$, $B$ bij $Q$ en $C$ bij $R$.
\end{stappen}
Dus is $\hoek{Q} = \hoek{B}$.}

Welk kenmerk gebruik je? Dat hangt af van wat je weet.

\begin{werkwijze}
Ken je hoeken, zoek dan twee paren gelijke hoeken (HH). Denk aan de hoekensom, aan overstaande hoeken, aan hoeken bij evenwijdige rechten en aan een hoek die twee driehoeken gemeenschappelijk hebben. Ken je zijden, rangschik ze dan en vergelijk de verhoudingen: van alle drie de zijden (ZZZ), of van twee zijden met de hoek ertussen (ZHZ).
\end{werkwijze}

\begin{oefeningen}
\oef{h2-zhz}
\end{oefeningen}

%========================================================================================
\section{Meetkundige problemen}

\subsection{Een evenwijdige in een driehoek}

Trek in een driehoek een rechte evenwijdig met een zijde. Dan snijdt ze een kleine driehoek af. Waarom heeft die kleine driehoek dezelfde vorm als de grote? De hoek in $A$ hebben ze gemeenschappelijk. En de hoeken in $D$ en $B$ zijn F-hoeken bij de evenwijdige rechten $DE$ en $BC$, dus gelijk. Twee gelijke hoeken: kenmerk HH.

\begin{theorie}
\begin{minipage}[t]{.62\linewidth}
\vspace{0pt}
\begin{stelling}{Evenwijdige met een zijde}
Ligt $D$ op $[AB]$ en $E$ op $[AC]$, met $DE \parallel BC$, dan is $\drieh{ADE} \gv \drieh{ABC}$, en dus

\vspace{1mm}
\qquad $\formule{\dfrac{\lengte{AD}}{\lengte{AB}} = \dfrac{\lengte{AE}}{\lengte{AC}} = \dfrac{\lengte{DE}}{\lengte{BC}}}$
\end{stelling}
\end{minipage}\hfill
\begin{minipage}[t]{.34\linewidth}
\vspace{0pt}
\centering
\begin{tikzpicture}[scale=.8,font=\scriptsize]
  \coordinate (A) at (1.2,2.6); \coordinate (B) at (0,0); \coordinate (C) at (3.6,0);
  \coordinate (D) at ($(A)!.45!(B)$); \coordinate (E) at ($(A)!.45!(C)$);
  \filldraw[fill=accent!10,draw=accent,thick] (A) -- (B) -- (C) -- cycle;
  \filldraw[fill=accent!30,draw=accent,thick] (A) -- (D) -- (E) -- cycle;
  \node[above] at (A) {$A$}; \node[below] at (B) {$B$}; \node[below] at (C) {$C$};
  \node[left] at (D) {$D$}; \node[right] at (E) {$E$};
\end{tikzpicture}
\end{minipage}
\end{theorie}

\alleenvoor{1,2,3,4}{Deze eigenschap heet ook de \emph{stelling van Thales}, naar de Griekse wiskundige die er volgens de overlevering de hoogte van een piramide mee bepaalde.}\alleenvoor{5}{Ze is een bijzonder geval van de \emph{stelling van Thales} hieronder, genoemd naar de Griekse wiskundige die er volgens de overlevering de hoogte van een piramide mee bepaalde.} Let op de valkuil: de factor is de verhouding van $\lengte{AD}$ tot de hele zijde $\lengte{AB}$, niet tot het stuk $\lengte{DB}$.

\voorbeeld{}{In driehoek $ABC$ is $DE \parallel BC$, met $\lengte{AD} = 3$, $\lengte{DB} = 2$ en $\lengte{DE} = 4{,}5$. Bereken $\lengte{BC}$.\label{vb:h2-ev}}{%
\uitleg{De kleine driehoek $ADE$ en de grote $ABC$ zijn gelijkvormig. De zijde $[AD]$ hoort bij de hele zijde $[AB]$, en die is $3 + 2 = 5$ lang.}
\naast{\begin{stappenvrij}{5.4cm}{4.4cm}
\lengte{AB} = 3 + 2 = 5 & De hele zijde. \\[1ex]
k = \dfrac{\lengte{AB}}{\lengte{AD}} = \dfrac53 & Van de kleine naar de grote driehoek. \\[2ex]
\lengte{BC} = \dfrac53 \cdot 4{,}5 = 7{,}5 & $[BC]$ hoort bij $[DE]$.
\end{stappenvrij}}{\begin{tikzpicture}[scale=.62,font=\scriptsize]
  \coordinate (A) at (1.5,3.5); \coordinate (B) at (0,0); \coordinate (C) at (5,0);
  \coordinate (D) at ($(A)!.6!(B)$); \coordinate (E) at ($(A)!.6!(C)$);
  \filldraw[fill=accent!10,draw=accent,thick] (A) -- (B) -- (C) -- cycle;
  \filldraw[fill=accent!30,draw=accent,thick] (A) -- (D) -- (E) -- cycle;
  \node[above] at (A) {$A$}; \node[below] at (B) {$B$}; \node[below] at (C) {$C$};
  \node[left] at (D) {$D$}; \node[right] at (E) {$E$};
  \node[left=1pt] at ($(A)!.5!(D)$) {3}; \node[left=1pt] at ($(D)!.5!(B)$) {2};
  \node[above] at ($(D)!.5!(E)$) {4,5}; \node[below] at ($(B)!.5!(C)$) {?};
\end{tikzpicture}}
Wie de fout maakt en met $\tfrac32$ rekent, vindt 6,75. Die vergelijkt $[DB]$ met $[AD]$, maar $[DB]$ is geen zijde van een van beide driehoeken.}

\voorbeeld{}{In driehoek $ABC$ is $DE \parallel BC$, met $\lengte{AD} = 4$, $\lengte{AB} = 10$ en $\lengte{AC} = 12$. Bereken $\lengte{AE}$ en $\lengte{EC}$.\label{vb:h2-ev2}}{%
\uitleg{Nu zoeken we een zijde van de kleine driehoek. We vermenigvuldigen met de factor van groot naar klein.}
\begin{stappen}
k = \dfrac{\lengte{AD}}{\lengte{AB}} = \dfrac{4}{10} = 0{,}4 & Van de grote naar de kleine driehoek. \\[2ex]
\lengte{AE} = 0{,}4 \cdot 12 = 4{,}8 & $[AE]$ hoort bij $[AC]$. \\[1ex]
\lengte{EC} = 12 - 4{,}8 = 7{,}2 & De rest van de zijde.
\end{stappen}}

\begin{oefeningen}
\oef{h2-ev}
\end{oefeningen}

\begin{alleenvoor}{5}
\subsection{De stelling van Thales}

Een ladder staat schuin tegen een muur, en de zon schijnt erop. Alle zonnestralen lopen evenwijdig, dus elk punt van de ladder heeft een schaduwpunt op de grond: schuif het punt langs een zonnestraal tot op de grond. De sporten van de ladder liggen even ver uit elkaar. Liggen hun schaduwen dan ook even ver uit elkaar? De ladder en de grond hoeven daarvoor niet door één punt te gaan, en dan helpt de evenwijdige met een zijde niet meteen.

\begin{theorie}
Kies een rechte $b$ en een rechte $r$ die $b$ snijdt. De \textbf{\textcolor{accent}{evenwijdige projectie}} op $b$ volgens $r$ beeldt een punt $P$ af op het snijpunt $P'$ van $b$ met de rechte door $P$ evenwijdig met $r$. Staat $r$ loodrecht op $b$, dan heet ze een \textbf{\textcolor{accent}{loodrechte projectie}}.
\end{theorie}

De schaduw van de ladder is zo'n projectie: $b$ is de grond, en $r$ is de richting van de zonnestralen. Wat gebeurt er met de afstanden tussen de punten?

\begin{theorie}
\begin{minipage}[t]{.56\linewidth}
\vspace{0pt}
\begin{stelling}{Stelling van Thales}
Liggen $A$, $B$ en $C$ op een rechte $a$, en zijn $A'$, $B'$ en $C'$ hun evenwijdige projecties op een rechte $b$, dan is

\vspace{1mm}
\qquad $\formule{\dfrac{\lengte{A'B'}}{\lengte{B'C'}} = \dfrac{\lengte{AB}}{\lengte{BC}}}$
\end{stelling}

\begin{bewijs}
Neem $B$ tussen $A$ en $C$. Is $a \parallel b$, dan zijn $ABB'A'$ en $BCC'B'$ parallellogrammen, en zijn de stukken gewoon even lang. Anders trek je door $A'$ de evenwijdige met $a$. Ze snijdt $BB'$ in $B_1$ en $CC'$ in $C_1$. De vierhoeken $AA'B_1B$ en $BB_1C_1C$ zijn parallellogrammen, dus $\lengte{A'B_1} = \lengte{AB}$ en $\lengte{B_1C_1} = \lengte{BC}$, zodat ook $\lengte{A'C_1} = \lengte{AC}$. In driehoek $A'C_1C'$ is $B_1B' \parallel C_1C'$. De evenwijdige met een zijde geeft

\vspace{1mm}
\qquad $\dfrac{\lengte{A'C'}}{\lengte{A'B'}} = \dfrac{\lengte{A'C_1}}{\lengte{A'B_1}} = \dfrac{\lengte{AC}}{\lengte{AB}}$

\vspace{1mm}
Trek van beide kanten 1 af: $\dfrac{\lengte{B'C'}}{\lengte{A'B'}} = \dfrac{\lengte{BC}}{\lengte{AB}}$, en neem het omgekeerde.
\end{bewijs}
\end{minipage}\hfill
\begin{minipage}[t]{.4\linewidth}
\vspace{0pt}
\centering
\begin{tikzpicture}[scale=.66,font=\scriptsize]
  \coordinate (A2) at (0,0); \coordinate (B2) at (2,0); \coordinate (C2) at (5,0);
  \coordinate (A) at (.868,1.578); \coordinate (B) at (3.295,2.354); \coordinate (C) at (6.936,3.519);
  \coordinate (B1) at (2.427,.777); \coordinate (C1) at (6.068,1.942);
  \draw[gegeven] (-.5,0) -- (6.2,0) node[below] {$b$};
  \draw[gegeven] ($(A)!-.25!(C)$) -- ($(A)!1.08!(C)$) node[above] {$a$};
  \draw[accent,thick] ($(A2)!-.12!(A)$) -- ($(A2)!1.2!(A)$);
  \draw[accent,thick] ($(B2)!-.1!(B)$) -- ($(B2)!1.12!(B)$);
  \draw[accent,thick] ($(C2)!-.08!(C)$) -- ($(C2)!1.08!(C)$) node[above] {$r$};
  \draw[grijs,dashed] (A2) -- ($(A2)!1.08!(C1)$);
  \foreach \p in {A,B,C,A2,B2,C2,B1,C1} \fill (\p) circle (1.3pt);
  \node[above left] at (A) {$A$}; \node[above left] at (B) {$B$}; \node[above left] at (C) {$C$};
  \node[below] at (A2) {$A'$}; \node[below] at (B2) {$B'$}; \node[below] at (C2) {$C'$};
  \node[right] at (B1) {$B_1$}; \node[right] at (C1) {$C_1$};
\end{tikzpicture}
\end{minipage}
\end{theorie}

In woorden: evenwijdige rechten snijden op twee rechten stukken af die evenredig zijn. Een evenwijdige projectie behoudt dus de verhouding van lijnstukken op een rechte, al verandert ze de lengtes zelf. De schaduwen van de sporten liggen dus ook even ver uit elkaar, en het midden van een lijnstuk wordt het midden van zijn projectie. De evenwijdige met een zijde is het geval waarin $a$ en $b$ door hetzelfde punt $A$ gaan: $A$ is dan zijn eigen projectie.

\voorbeeld{}{Drie evenwijdige straten snijden de Molenlaan in $A$, $B$ en $C$, en de Kerklaan in $A'$, $B'$ en $C'$. Langs de Molenlaan is $\lengte{AB} = 120$~m en $\lengte{BC} = 180$~m. Langs de Kerklaan is $\lengte{A'B'} = 150$~m. Hoe lang is $[B'C']$?\label{vb:h2-thales}}{%
\uitleg{De punten op de Kerklaan zijn de evenwijdige projecties van de punten op de Molenlaan, volgens de richting van de straten. De stelling van Thales geeft de verhouding.}
\naast{\begin{stappenvrij}{5.2cm}{3.4cm}
\dfrac{\lengte{B'C'}}{\lengte{A'B'}} = \dfrac{\lengte{BC}}{\lengte{AB}} = \dfrac{180}{120} = 1{,}5 & Stelling van Thales. \\[2ex]
\lengte{B'C'} = 1{,}5 \cdot 150 = 225~\text{m} & Het gevraagde stuk.
\end{stappenvrij}}{\begin{tikzpicture}[scale=.44,font=\scriptsize]
  \coordinate (A2) at (0,0); \coordinate (B2) at (2.5,0); \coordinate (C2) at (6.25,0);
  \coordinate (A) at (.949,2.109); \coordinate (B) at (3.724,2.719); \coordinate (C) at (7.886,3.635);
  \draw[gegeven] (7,0) -- (-.6,0) node[left] {Kerklaan};
  \draw[gegeven] ($(A)!1.08!(C)$) -- ($(A)!-.12!(C)$) node[left] {Molenlaan};
  \foreach \p/\q in {A2/A,B2/B,C2/C} \draw[accent,thick] ($(\p)!-.2!(\q)$) -- ($(\p)!1.25!(\q)$);
  \foreach \p in {A,B,C,A2,B2,C2} \fill (\p) circle (1.5pt);
  \node[above left] at (A) {$A$}; \node[above left] at (B) {$B$}; \node[above left] at (C) {$C$};
  \node[below] at (A2) {$A'$}; \node[below] at (B2) {$B'$}; \node[below] at (C2) {$C'$};
  \node[below] at ($(A2)!.5!(B2)$) {150}; \node[below] at ($(B2)!.5!(C2)$) {?};
  \node[above] at ($(A)!.5!(B)$) {120}; \node[above] at ($(B)!.5!(C)$) {180};
\end{tikzpicture}}
Langs de Kerklaan liggen de kruispunten verder uit elkaar dan langs de Molenlaan, maar in dezelfde verhouding.}

\begin{werkwijze}
Zoek de evenwijdige rechten, en de twee rechten die ze snijden. Schrijf de verhouding van twee stukken op de ene rechte gelijk aan de verhouding van de overeenkomstige stukken op de andere, en los op.
\end{werkwijze}

\begin{oefeningen}
\oef{h2-thales}
\end{oefeningen}
\end{alleenvoor}

\subsection{De zandloper}

Snijden twee rechten elkaar in een punt $S$, en verbind je ze aan beide kanten met evenwijdige lijnstukken, dan krijg je een figuur die op een zandloper lijkt. Ook daar zitten twee gelijkvormige driehoeken in. De hoeken in $S$ zijn overstaande hoeken, dus gelijk. En de hoeken in $A$ en $C$ zijn Z-hoeken bij de evenwijdige rechten $AB$ en $CD$.\alleenvoor{4,5}{ De ene driehoek is het beeld van de andere bij een homothetie met centrum $S$ en een negatieve factor.}

\voorbeeld{}{De lijnstukken $[AC]$ en $[BD]$ snijden elkaar in $S$, en $AB \parallel CD$. Verder is $\lengte{SA} = 3$, $\lengte{SC} = 5$ en $\lengte{AB} = 4{,}2$. Bereken $\lengte{CD}$.\label{vb:h2-zand}}{%
\uitleg{Eerst tonen we aan dat de driehoeken gelijkvormig zijn, en we letten op de volgorde: $S$ bij $S$, $A$ bij $C$ en $B$ bij $D$.}
\naast{\begin{stappenvrij}{5.2cm}{4.6cm}
\hoek{ASB} = \hoek{CSD} & Overstaande hoeken. \\[1ex]
\hoek{A} = \hoek{C} & Z-hoeken, want $AB \parallel CD$. \\[1ex]
\drieh{SAB} \gv \drieh{SCD} & Kenmerk HH. \\[1ex]
k = \dfrac{\lengte{SC}}{\lengte{SA}} = \dfrac53 & $[SC]$ hoort bij $[SA]$. \\[2ex]
\lengte{CD} = \dfrac53 \cdot 4{,}2 = 7 & $[CD]$ hoort bij $[AB]$.
\end{stappenvrij}}{\begin{tikzpicture}[scale=.5,font=\scriptsize]
  \coordinate (S) at (0,0);
  \coordinate (A) at (-1.2,1.8); \coordinate (B) at (1.32,1.8);
  \coordinate (C) at (2,-3); \coordinate (D) at (-2.2,-3);
  \filldraw[fill=accent!30,draw=accent,thick] (S) -- (A) -- (B) -- cycle;
  \filldraw[fill=accent!12,draw=accent,thick] (S) -- (C) -- (D) -- cycle;
  \hoekboog{S}{A}{B}{3mm}
  \hoekboog{S}{C}{D}{3mm}
  \node[above] at (A) {$A$}; \node[above] at (B) {$B$}; \node[below] at (C) {$C$}; \node[below] at (D) {$D$};
  \node[right=2pt] at (S) {$S$};
\end{tikzpicture}}}

\begin{oefeningen}
\oef{h2-zand}
\end{oefeningen}

\begin{alleenvoor}{5}
\subsection{Middenparallel en zwaartelijnen}

Verbind je in een driehoek de middens van twee zijden, dan krijg je een \emph{middenparallel}. Teken er een paar, en je vermoedt al wat de naam zegt: hij loopt evenwijdig met de derde zijde, en is half zo lang. Met de kenmerken kun je bewijzen dat dat altijd zo is.

\begin{theorie}
\begin{minipage}[t]{.62\linewidth}
\vspace{0pt}
\begin{stelling}{Middenparallel}
Is $M$ het midden van $[AB]$ en $N$ het midden van $[AC]$, dan is

\vspace{1mm}
\qquad $\formule{MN \parallel BC} \quad \text{en} \quad \formule{\lengte{MN} = \tfrac12\,\lengte{BC}}$
\end{stelling}

\begin{bewijs}
In $\drieh{AMN}$ en $\drieh{ABC}$ is $\hoek{A}$ gemeenschappelijk, en

\vspace{1mm}
\qquad $\dfrac{\lengte{AM}}{\lengte{AB}} = \dfrac{\lengte{AN}}{\lengte{AC}} = \dfrac12$

\vspace{1mm}
Volgens ZHZ is $\drieh{AMN} \gv \drieh{ABC}$ met factor $\tfrac12$. Dus is $\lengte{MN} = \tfrac12\,\lengte{BC}$, en $\hoek{AMN} = \hoek{ABC}$. Dat zijn F-hoeken bij de rechten $MN$ en $BC$ en de snijlijn $AB$, dus $MN \parallel BC$.
\end{bewijs}
\end{minipage}\hfill
\begin{minipage}[t]{.34\linewidth}
\vspace{0pt}
\centering
\begin{tikzpicture}[scale=.85,font=\scriptsize]
  \coordinate (A) at (1.4,3); \coordinate (B) at (0,0); \coordinate (C) at (4.4,0);
  \coordinate (M) at ($(A)!.5!(B)$); \coordinate (N) at ($(A)!.5!(C)$);
  \filldraw[fill=accent!10,draw=accent,thick] (A) -- (B) -- (C) -- cycle;
  \filldraw[fill=accent!30,draw=accent,thick] (A) -- (M) -- (N) -- cycle;
  \coordinate (S1) at ($(A)!.5!(M)$); \coordinate (S2) at ($(M)!.5!(B)$);
  \coordinate (S3) at ($(A)!.44!(N)$); \coordinate (S4) at ($(A)!.56!(N)$);
  \coordinate (S5) at ($(N)!.44!(C)$); \coordinate (S6) at ($(N)!.56!(C)$);
  \streep{S1}{A} \streep{S2}{B} \streep{S3}{A} \streep{S4}{A} \streep{S5}{C} \streep{S6}{C}
  \node[above] at (A) {$A$}; \node[below] at (B) {$B$}; \node[below] at (C) {$C$};
  \node[left] at (M) {$M$}; \node[right] at (N) {$N$};
\end{tikzpicture}
\end{minipage}
\end{theorie}

Zo heeft een driehoek drie middenparallellen. Samen vormen ze een kleine driehoek in het midden, de \emph{middendriehoek}.

\voorbeeld{}{Een driehoek heeft zijden van 8, 10 en 12~cm en een oppervlakte van ongeveer 39,7~cm$^2$. Hoe groot zijn de omtrek en de oppervlakte van zijn middendriehoek?\label{vb:h2-midden}}{%
\uitleg{Elke zijde van de middendriehoek is een middenparallel, dus half zo lang als een zijde van de grote driehoek. De middendriehoek is dus gelijkvormig met de grote, met factor $\tfrac12$ (ZZZ).}
\begin{stappen}
4, \ 5 \text{ en } 6~\text{cm} & De zijden van de middendriehoek. \\[1ex]
4 + 5 + 6 = 15~\text{cm} & De omtrek: de helft van $8 + 10 + 12 = 30$. \\[1ex]
\left(\tfrac12\right)^2 \cdot 39{,}7 \approx 9{,}9~\text{cm}^2 & Een oppervlakte wordt met $k^2$ vermenigvuldigd.
\end{stappen}
De drie middenparallellen verdelen de driehoek in vier gelijke driehoekjes: elk een vierde van de oppervlakte.}

Een \emph{zwaartelijn} verbindt een hoekpunt met het midden van de overstaande zijde. Teken je de drie zwaartelijnen, dan lijken ze door één punt te gaan. Waarom, en waar ligt dat punt? Een middenparallel en een zandloper geven het antwoord.

\begin{theorie}
\begin{minipage}[t]{.62\linewidth}
\vspace{0pt}
\begin{stelling}{Zwaartelijnen}
De drie zwaartelijnen van een driehoek gaan door één punt $Z$, het \textbf{\textcolor{accent}{zwaartepunt}}. Is $M$ het midden van $[BC]$, dan is

\vspace{1mm}
\qquad $\formule{\lengte{AZ} = 2\,\lengte{ZM}}$

\vspace{1mm}
en zo ook voor de andere zwaartelijnen.
\end{stelling}

\begin{bewijs}
Neem ook het midden $N$ van $[AC]$, en noem $Z$ het snijpunt van de zwaartelijnen $[AM]$ en $[BN]$. In $\drieh{CBA}$ is $MN$ een middenparallel: $MN \parallel AB$ en $\lengte{MN} = \tfrac12\,\lengte{AB}$. In de zandloper zijn de hoeken in $Z$ overstaand, en $\hoek{ZAB} = \hoek{ZMN}$ (Z-hoeken). Volgens HH is $\drieh{ZAB} \gv \drieh{ZMN}$ met factor $\tfrac12$, dus $\lengte{ZM} = \tfrac12\,\lengte{ZA}$.

Zo ligt $Z$ op $[AM]$, twee keer zo ver van $A$ als van $M$. Neem je de derde zwaartelijn in plaats van $[BN]$, dan vind je met dezelfde redenering hetzelfde punt van $[AM]$. De drie zwaartelijnen gaan dus door $Z$.
\end{bewijs}
\end{minipage}\hfill
\begin{minipage}[t]{.34\linewidth}
\vspace{0pt}
\centering
\begin{tikzpicture}[scale=.72,font=\scriptsize]
  \coordinate (A) at (0,0); \coordinate (B) at (6,0); \coordinate (C) at (2,4);
  \coordinate (M) at (4,2); \coordinate (N) at (1,2); \coordinate (Z) at (2.667,1.333);
  \filldraw[fill=accent!10,draw=accent,thick] (A) -- (B) -- (C) -- cycle;
  \fill[accent!30] (Z) -- (A) -- (B) -- cycle; \fill[accent!30] (Z) -- (M) -- (N) -- cycle;
  \draw[accent,thick] (A) -- (M) (B) -- (N);
  \draw[accent,dashed] (M) -- (N);
  \fill (Z) circle (1.5pt);
  \node[below] at (A) {$A$}; \node[below] at (B) {$B$}; \node[above] at (C) {$C$};
  \node[right] at (M) {$M$}; \node[left] at (N) {$N$}; \node[below=2pt] at (Z) {$Z$};
\end{tikzpicture}
\end{minipage}
\end{theorie}

Het zwaartepunt verdeelt elke zwaartelijn dus in een stuk van twee derde (aan de kant van het hoekpunt) en een stuk van een derde (aan de kant van het midden). Knip je een driehoek uit karton, dan blijft hij op een punt van een potlood in evenwicht als je hem in het zwaartepunt ondersteunt: vandaar de naam.

\voorbeeld{}{De zwaartelijnen $[AM]$ en $[BN]$ van driehoek $ABC$ snijden elkaar in $Z$. Verder is $\lengte{AM} = 12$ en $\lengte{BZ} = 10$. Bereken $\lengte{AZ}$, $\lengte{ZM}$, $\lengte{ZN}$ en $\lengte{BN}$.\label{vb:h2-zwaarte}}{%
\uitleg{$Z$ is het zwaartepunt. Het stuk naar het hoekpunt is twee derde van de zwaartelijn, het stuk naar het midden een derde, of de helft van het stuk naar het hoekpunt.}
\begin{stappen}
\lengte{AZ} = \tfrac23 \cdot 12 = 8, \quad \lengte{ZM} = \tfrac13 \cdot 12 = 4 & Op de zwaartelijn $[AM]$. \\[1ex]
\lengte{ZN} = \tfrac12 \cdot \lengte{BZ} = 5 & Op de zwaartelijn $[BN]$ is $\lengte{BZ} = 2\,\lengte{ZN}$. \\[1ex]
\lengte{BN} = 10 + 5 = 15 & De hele zwaartelijn.
\end{stappen}}

\begin{oefeningen}
\oef{h2-midden}
\end{oefeningen}
\end{alleenvoor}

\subsection{Hoogtes en afstanden}

Hoe meet je de hoogte van een boom zonder erin te klimmen? Met gelijkvormige driehoeken: je maakt een kleine driehoek die je wel kunt meten, met dezelfde hoeken als een grote driehoek die je niet kunt meten. Een spiegel is daarvoor handig, omdat licht weerkaatst onder dezelfde hoek als waaronder het invalt.

\voorbeeld{}{Jana legt een spiegeltje op de grond, 6~m van een boom. Ze stapt achteruit tot ze de top van de boom in de spiegel ziet. Ze staat dan 1,5~m van de spiegel, en haar ogen zijn 1,6~m boven de grond. Hoe hoog is de boom?\label{vb:h2-spiegel}}{%
\uitleg{Er zijn twee rechthoekige driehoeken: Jana met de spiegel, en de boom met de spiegel. In de spiegel zijn de invalshoek en de terugkaatsingshoek gelijk.}
\naast{\begin{stappenvrij}{5.2cm}{4cm}
\hoek{J} = \hoek{V} = 90^\circ & Jana en de boom staan recht. \\[1ex]
\hoek{OSJ} = \hoek{TSV} & Weerkaatsing in de spiegel. \\[1ex]
\drieh{SJO} \gv \drieh{SVT} & Kenmerk HH. \\[1ex]
\dfrac{\lengte{VT}}{\lengte{JO}} = \dfrac{\lengte{SV}}{\lengte{SJ}} \ \Rightarrow \ \dfrac{h}{1{,}6} = \dfrac{6}{1{,}5} = 4 & De overeenkomstige zijden. \\[2ex]
h = 4 \cdot 1{,}6 = 6{,}4~\text{m} & De hoogte van de boom.
\end{stappenvrij}}{\begin{tikzpicture}[scale=.5,font=\scriptsize]
  \draw[lijngrijs] (-2.3,0) -- (6.6,0);
  \coordinate (J) at (-1.5,0); \coordinate (O) at (-1.5,1.6); \coordinate (S) at (0,0);
  \coordinate (V) at (6,0); \coordinate (T) at (6,6.4);
  \draw[accent!60!black,line width=2pt] (6,0) -- (6,4.4);
  \fill[accent!25] (6,5.2) circle (1.3);
  \draw[accent,thick] (O) -- (S) -- (T);
  \draw[grijs,thick] (J) -- (O) (V) -- (T);
  \fill (O) circle (1.5pt); \fill (T) circle (1.5pt);
  \fill[grijs] (-.3,-.06) rectangle (.3,.06);
  \hoekboog{V}{S}{T}{3mm} \hoekboog{O}{S}{J}{3mm}
  \node[below] at (J) {$J$}; \node[left] at (O) {$O$}; \node[below] at (S) {$S$};
  \node[below] at (V) {$V$}; \node[right] at (T) {$T$};
\end{tikzpicture}}}

\begin{oefeningen}
\oef{h2-afstand}
\end{oefeningen}

Nu kunnen we het probleem oplossen. Het idee is hetzelfde: twee driehoeken met dezelfde hoeken, een kleine in de doos en een grote buiten.

\voorbeeld{}{Terug naar het probleem. Een kerktoren staat 150~m voor het gaatje van een camera. Het scherm staat 30~cm achter het gaatje, en het beeld van de toren is daar 7~cm hoog. Hoe hoog is de toren?\label{vb:h2-terug}}{%
\uitleg{We nemen aan dat het gaatje $G$ even hoog staat als de voet $V$ van de toren, en dat het scherm rechtop staat. Het licht van de top $T$ gaat in een rechte lijn door het gaatje naar het beeld $T'$ op het scherm, en $V'$ ligt recht achter het gaatje.}
\naast{\begin{stappenvrij}{5.2cm}{4.4cm}
\hoek{V} = \hoek{V'} = 90^\circ & De toren en het scherm staan rechtop. \\[1ex]
\hoek{VGT} = \hoek{V'GT'} & Overstaande hoeken in het gaatje. \\[1ex]
\drieh{GVT} \gv \drieh{GV'T'} & Kenmerk HH: een zandloper. \\[1ex]
\dfrac{\lengte{VT}}{\lengte{V'T'}} = \dfrac{\lengte{GV}}{\lengte{GV'}} = \dfrac{15\,000}{30} = 500 & Alles in centimeter. \\[2ex]
\lengte{VT} = 500 \cdot 7 = 3500~\text{cm} & De hoogte van de toren.
\end{stappenvrij}}{\begin{tikzpicture}[scale=.55,font=\scriptsize]
  \coordinate (V) at (0,0); \coordinate (T) at (0,3.5); \coordinate (G) at (5,0);
  \coordinate (V2) at (6.2,0); \coordinate (T2) at (6.2,-.84);
  \draw[accent,thick] (V) -- (T) -- (T2) -- (V2) -- cycle;
  \fill[accent!30] (V) -- (T) -- (G) -- cycle;
  \fill[accent!30] (G) -- (V2) -- (T2) -- cycle;
  \rechtehoek{T}{V}{G} \rechtehoek{G}{V2}{T2}
  \node[below left] at (V) {$V$}; \node[left] at (T) {$T$}; \node[below] at (G) {$G$};
  \node[above] at (V2) {$V'$}; \node[below] at (T2) {$T'$};
\end{tikzpicture}}
De toren is 35~m hoog. De figuur is niet op schaal: in werkelijkheid is de grote driehoek 500 keer zo groot als de kleine.}

Zo is het probleem opgelost. Een verhouding opstellen mocht, omdat de twee driehoeken gelijkvormig zijn, en dat volgt uit twee gelijke hoeken: een rechte hoek en overstaande hoeken. Zonder die redenering is een verhouding maar een gok.

\begin{werkwijze}
Maak een figuur en zoek twee driehoeken: een met de lengte die je zoekt, een andere met lengtes die je kent. Toon met een kenmerk aan dat ze gelijkvormig zijn, meestal met HH. Schrijf de overeenkomstige hoekpunten in dezelfde volgorde, stel de verhouding op en los op. Kijk of je antwoord kan kloppen.
\end{werkwijze}

\begin{alleenvoor}{4,5}
\subsection{De hoogtelijn in een rechthoekige driehoek}

Een rechthoekige driehoek valt door zijn hoogtelijn uiteen in twee kleinere rechthoekige driehoeken. Alle drie hebben ze dezelfde hoeken, en dat levert een handige verhouding op.

\voorbeeld{}{Driehoek $ABC$ is rechthoekig in $A$, en $[AD]$ is de hoogtelijn op $[BC]$. Verder is $\lengte{BD} = 4$ en $\lengte{DC} = 9$. Bereken $\lengte{AD}$.\label{vb:h2-hoogte}}{%
\uitleg{De driehoeken $DBA$ en $DAC$ zijn allebei rechthoekig in $D$. Daarnaast is $\hoek{B} = 90^\circ - \hoek{C}$, en in driehoek $DAC$ is ook $\hoek{DAC} = 90^\circ - \hoek{C}$.}
\naast{\begin{stappenvrij}{5.2cm}{4.4cm}
\hoek{BDA} = \hoek{ADC} = 90^\circ & $[AD]$ is een hoogtelijn. \\[1ex]
\hoek{B} = \hoek{DAC} & Beide zijn $90^\circ - \hoek{C}$. \\[1ex]
\drieh{DBA} \gv \drieh{DAC} & Kenmerk HH. \\[1ex]
\dfrac{\lengte{AD}}{\lengte{BD}} = \dfrac{\lengte{DC}}{\lengte{AD}} & Overeenkomstige zijden. \\[2ex]
\lengte{AD}^2 = 4 \cdot 9 = 36 & Kruiselings vermenigvuldigen. \\[1ex]
\lengte{AD} = 6 & Een lengte is positief.
\end{stappenvrij}}{\begin{tikzpicture}[scale=.32,font=\scriptsize]
  \coordinate (B) at (0,0); \coordinate (C) at (13,0); \coordinate (D) at (4,0); \coordinate (A) at (4,6);
  \filldraw[fill=accent!30,draw=accent,thick] (B) -- (D) -- (A) -- cycle;
  \filldraw[fill=accent!12,draw=accent,thick] (D) -- (C) -- (A) -- cycle;
  \rechtehoek{A}{D}{C} \rechtehoek{B}{A}{C}
  \node[below] at (B) {$B$}; \node[below] at (C) {$C$}; \node[below] at (D) {$D$}; \node[above] at (A) {$A$};
  \node[below] at (2,0) {4}; \node[below] at (8.5,0) {9};
\end{tikzpicture}}}

\begin{alleenvoor}{5}
Werkt dat in elke rechthoekige driehoek? De hoogtelijn $[AD]$ verdeelt de schuine zijde in twee stukken: $[BD]$ is de loodrechte projectie van de rechthoekszijde $[AB]$ op de schuine zijde, en $[DC]$ die van $[AC]$. Gelijkheden tussen zulke lengtes heten \emph{metrische betrekkingen}. Er zijn er twee, en je bewijst ze allebei met gelijkvormige driehoeken.

\begin{theorie}
\begin{minipage}[t]{.62\linewidth}
\vspace{0pt}
\begin{stelling}{Hoogte in een rechthoekige driehoek}
Is $\drieh{ABC}$ rechthoekig in $A$, en is $[AD]$ de hoogtelijn op $[BC]$, dan is

\vspace{1mm}
\qquad $\formule{\lengte{AD}^2 = \lengte{BD} \cdot \lengte{DC}}$
\end{stelling}

\begin{bewijs}
De driehoeken $DBA$ en $DAC$ zijn rechthoekig in $D$, en $\hoek{B} = 90^\circ - \hoek{C} = \hoek{DAC}$. Volgens HH is $\drieh{DBA} \gv \drieh{DAC}$, dus

\vspace{1mm}
\qquad $\dfrac{\lengte{AD}}{\lengte{BD}} = \dfrac{\lengte{DC}}{\lengte{AD}}$, \quad en kruiselings $\lengte{AD}^2 = \lengte{BD} \cdot \lengte{DC}$
\end{bewijs}
\end{minipage}\hfill
\begin{minipage}[t]{.34\linewidth}
\vspace{0pt}
\centering
\begin{tikzpicture}[scale=.8,font=\scriptsize]
  \coordinate (B) at (0,0); \coordinate (C) at (5,0); \coordinate (D) at (1.8,0); \coordinate (A) at (1.8,2.4);
  \filldraw[fill=accent!30,draw=accent,thick] (B) -- (D) -- (A) -- cycle;
  \filldraw[fill=accent!12,draw=accent,thick] (D) -- (C) -- (A) -- cycle;
  \rechtehoek{A}{D}{C} \rechtehoek{B}{A}{C}
  \node[below] at (B) {$B$}; \node[below] at (C) {$C$}; \node[below] at (D) {$D$}; \node[above] at (A) {$A$};
\end{tikzpicture}
\end{minipage}
\end{theorie}

In woorden: het kwadraat van de hoogte is het product van de twee stukken van de schuine zijde. Voor de rechthoekszijden vergelijk je niet de twee kleine driehoeken met elkaar, maar een kleine driehoek met de grote.

\begin{theorie}
\begin{minipage}[t]{.62\linewidth}
\vspace{0pt}
\begin{stelling}{Rechthoekszijde in een rechthoekige driehoek}
Is $\drieh{ABC}$ rechthoekig in $A$, en is $[AD]$ de hoogtelijn op $[BC]$, dan is

\vspace{1mm}
\qquad $\formule{\lengte{AB}^2 = \lengte{BD} \cdot \lengte{BC}}$ \quad en

\vspace{1mm}
\qquad $\formule{\lengte{AC}^2 = \lengte{CD} \cdot \lengte{CB}}$
\end{stelling}

\begin{bewijs}
De driehoeken $ABC$ en $DBA$ hebben een rechte hoek (in $A$ en in $D$), en de hoek $\hoek{B}$ gemeenschappelijk. Volgens HH is $\drieh{ABC} \gv \drieh{DBA}$, dus

\vspace{1mm}
\qquad $\dfrac{\lengte{AB}}{\lengte{DB}} = \dfrac{\lengte{BC}}{\lengte{BA}}$, \quad en kruiselings $\lengte{AB}^2 = \lengte{BD} \cdot \lengte{BC}$

\vspace{1mm}
Voor $[AC]$ gaat het net zo, met $\drieh{ABC} \gv \drieh{DAC}$ en de hoek $\hoek{C}$.
\end{bewijs}
\end{minipage}\hfill
\begin{minipage}[t]{.34\linewidth}
\vspace{0pt}
\centering
\begin{tikzpicture}[scale=.8,font=\scriptsize]
  \coordinate (B) at (0,0); \coordinate (C) at (5,0); \coordinate (D) at (1.8,0); \coordinate (A) at (1.8,2.4);
  \filldraw[fill=accent!10,draw=accent,thick] (A) -- (B) -- (C) -- cycle;
  \filldraw[fill=accent!35,draw=accent,thick] (D) -- (B) -- (A) -- cycle;
  \rechtehoek{A}{D}{C} \rechtehoek{B}{A}{C}
  \node[below] at (B) {$B$}; \node[below] at (C) {$C$}; \node[below] at (D) {$D$}; \node[above] at (A) {$A$};
\end{tikzpicture}
\end{minipage}
\end{theorie}

In woorden: het kwadraat van een rechthoekszijde is het product van haar projectie op de schuine zijde en de hele schuine zijde. Tel je de twee betrekkingen op, dan vind je een bekende stelling (opdracht~\ref{oef:h2-o31}).

\voorbeeld{}{Driehoek $ABC$ is rechthoekig in $A$, en $[AD]$ is de hoogtelijn op $[BC]$. Verder is $\lengte{BD} = 3{,}6$ en $\lengte{BC} = 10$. Bereken $\lengte{AB}$, $\lengte{AC}$ en $\lengte{AD}$.\label{vb:h2-metrisch}}{%
\uitleg{We kennen de schuine zijde en een stuk ervan. Met de betrekking voor de rechthoekszijde vinden we $\lengte{AB}$. Voor de andere lengtes hebben we eerst het tweede stuk nodig.}
\naast{\begin{stappenvrij}{6.4cm}{3.4cm}
\lengte{AB}^2 = 3{,}6 \cdot 10 = 36, \ \ \lengte{AB} = 6 & Rechthoekszijde: $[BD]$ is de projectie van $[AB]$. \\[1ex]
\lengte{DC} = 10 - 3{,}6 = 6{,}4 & Het tweede stuk. \\[1ex]
\lengte{AC}^2 = 6{,}4 \cdot 10 = 64, \ \ \lengte{AC} = 8 & Rechthoekszijde, met $[DC]$. \\[1ex]
\lengte{AD}^2 = 3{,}6 \cdot 6{,}4 = 23{,}04, \ \ \lengte{AD} = 4{,}8 & Hoogte: het product van de stukken.
\end{stappenvrij}}{\begin{tikzpicture}[scale=.3,font=\scriptsize]
  \coordinate (B) at (0,0); \coordinate (C) at (10,0); \coordinate (D) at (3.6,0); \coordinate (A) at (3.6,4.8);
  \filldraw[fill=accent!20,draw=accent,thick] (A) -- (B) -- (C) -- cycle;
  \draw[accent,thick] (A) -- (D);
  \rechtehoek{A}{D}{C} \rechtehoek{B}{A}{C}
  \node[below] at (B) {$B$}; \node[below] at (C) {$C$}; \node[below] at (D) {$D$}; \node[above] at (A) {$A$};
  \node[below] at (1.8,0) {3,6};
\end{tikzpicture}}
Controle met de oppervlakte: $\tfrac12 \cdot 6 \cdot 8 = 24$ en $\tfrac12 \cdot 10 \cdot 4{,}8 = 24$.}

\begin{werkwijze}
Zoek de rechte hoek, de hoogtelijn op de schuine zijde en de twee stukken waarin ze de schuine zijde verdeelt. Kies de betrekking met de lengte die je zoekt en twee lengtes die je kent. Lukt dat niet meteen, bereken dan eerst een ontbrekend stuk.
\end{werkwijze}

\needspace{6\baselineskip}
\begin{oefeningen}
\oef{h2-metrisch}
\end{oefeningen}
\end{alleenvoor}
\end{alleenvoor}

%========================================================================================
\opdrachtenkop

\begin{oefeningen}
\oef{h2-o1}

\oef{h2-o25}

\oef{h2-o2}

\oef{h2-o3}

\oef{h2-o4}

\oef{h2-o5}

\oef{h2-o26}

\oef{h2-o6}

\oef{h2-o7}

\oef{h2-o8}

\oef{h2-o9}

\oef{h2-o10}

\oef{h2-o11}

\oef{h2-o12}

\oef{h2-o27}

\oef{h2-o13}

\oef{h2-o14}

\oef{h2-o15}

\oef{h2-o16}

\oef{h2-o28}

\oef{h2-o17}

\oef{h2-o18}

\oef{h2-o30}

\oef{h2-o19}

\oef{h2-o20}

\oef{h2-o21}

\oef{h2-o29}

\oef{h2-o22}

\oef{h2-o31}

\oef{h2-o23}

\oef{h2-o24}
\end{oefeningen}

%----------------------------------------------------------------------------------------
\opdrachtsectie{Test jezelf}

\begin{oefeningen}
\oef{h2-t1}

\oef{h2-t2}

\oef{h2-t3}

\oef{h2-t4}

\oef{h2-t5}

\oef{h2-t6}

\oef{h2-t7}

\oef{h2-t8}

\oef{h2-t9}
\end{oefeningen}

\herhaalschema{%
  \fv{oef:h2-t1}{vb:h2-def}{oef:h2-def}%
  \fv{oef:h2-t2}{vb:h2-hh}{oef:h2-hh}%
  \fv{oef:h2-t3}{vb:h2-zzz}{oef:h2-zzz}%
  \fv{oef:h2-t4}{vb:h2-zhz}{oef:h2-zhz}%
  \fv{oef:h2-t5}{vb:h2-ev}{oef:h2-ev}%
  \fv{oef:h2-t6}{vb:h2-zand}{oef:h2-zand}%
  \fv{oef:h2-t7}{vb:h2-spiegel}{oef:h2-o7}%
  \alleenvoor{5}{\fv{oef:h2-t8}{vb:h2-thales}{oef:h2-thales}\fv{oef:h2-t9}{vb:h2-metrisch}{oef:h2-metrisch}}}

\ifzelfstudie
\section*{Maak je eigen samenvatting}
Sluit dit boek en schrijf in je schrift: wanneer twee driehoeken gelijkvormig zijn en hoe je dat noteert, hoe je overeenkomstige hoekpunten en zijden vindt, de drie kenmerken HH, ZZZ en ZHZ, de eigenschap van een evenwijdige met een zijde, de zandloper, en hoe je met gelijkvormige driehoeken een hoogte of afstand berekent.\alleenvoor{5}{ Schrijf ook op: wat een evenwijdige projectie is en wat de stelling van Thales zegt, de eigenschappen van de middenparallel en de zwaartelijnen, en de metrische betrekkingen in een rechthoekige driehoek, telkens met het idee van het bewijs.} Vergelijk daarna met de kaders.
\fi
